\section{Introduction}\label{s:introduction}

An exact linear equation may contain arbitrarily fine real information.
Indeed, \(d\) independent Gaussian equations determine a vector in
\(\R^d\) almost surely, because their coefficient matrix has full rank.
A streaming learner with finitely many persistent states cannot retain
the equations themselves at arbitrary precision. We study how many
fresh equations it needs to estimate a direction when the information
passed from one observation to the next is bounded in bits.

\begin{definition}[Finite-state Gaussian stream]\label{s:model}
Fix integers \(d\ge2\), \(M\ge0\), and \(T\ge0\), and an accuracy
\(0<\epsilon\le1/10\). For a signal \(s\in S^{d-1}\), the observations
are
\[
 X_t\sim N(0,I_d),\qquad Y_t=\ip{X_t}{s},\qquad 1\le t\le T,
\]
with the \(X_t\) independent. The learner reads these pairs in order.
At every sample index its persistent state has at most \(2^M\) possible
values. Its transition and stopping rules may depend on the index,
\(d,M,T,\epsilon\), the current state, the entire current pair, and
fresh randomness. Computation within a transition is unrestricted.
All information retained for a later sample must be in the finite state.

The learner stops at an index \(\tau\in\{0,\ldots,T\}\) and returns
\(\widehat s\in S^{d-1}\). Its output law depends only on the terminal
state, \(\tau\), and fresh randomness. Thus the last pair affects the
output only through the terminal state.

A shared random seed may choose the rules before the stream starts. The pair consisting of this seed
and the initial state is jointly independent of the signal and all
sample rows; in particular, conditional on the seed, initialization
carries no signal or sample information. Conditional on the seed,
fresh transition and output randomness has no additional signal or
sample information, and the rules have only the arguments just listed.

The original randomized experiment is required to be jointly
measurable in the signal, pre-generated rows, seed, initialization,
and fresh transition and output randomness. For almost every fixed
seed, pair-input probabilities may be Borel measurable, or measurable
in the completion of
\(\gamma_d(dx)\,dy\), where \(\gamma_d=N(0,I_d)\) and \(dy\) is
Lebesgue measure. Joint measurability of the original experiment is
required also in the latter case; separate measurability assertions
for individual seeds do not by themselves define its mixture.
The sample bound \(T\) is a deterministic finite horizon.
\end{definition}

The model charges persistent bits rather than real registers. It
permits a different rule at every index and does not limit the work
performed on the current sample.

\begin{theorem}[Explicit noiseless precision bound]\label{s:main}
Let \(M:\mathbb N\to\mathbb Z_{\ge0}\) satisfy \(M(d)=o(d^2)\).
There is a dimension threshold depending only on this memory sequence
such that, for every larger \(d\), every \(0<\epsilon\le1/10\), and
every learner in Definition~\ref{s:model}, the following implication
holds. If \(S\) is uniform on \(S^{d-1}\), independently of the
Gaussian rows and learner randomness, and
\[
 \Prob\{\arccos\ip{\widehat s}{S}\le\epsilon\}\ge\frac23,
\]
then
\begin{equation}\label{s:main-bound}
 T\ge 2^{-16}d\log_2(1/\epsilon).
\end{equation}
\end{theorem}

\begin{corollary}[Guarantees for every signal]\label{s:worst-signal}
Under the memory assumption of Theorem~\ref{s:main}, suppose a learner
satisfies
\[
 \Prob_s\{\arccos\ip{\widehat s}{s}\le\epsilon\}\ge\frac23
 \qquad\text{for every }s\in S^{d-1},
\]
where \(\Prob_s\) averages its observations and randomness at the fixed
signal \(s\). Then \eqref{s:main-bound} holds at the same dimension
threshold.
\end{corollary}
\begin{proof}
Integrate the displayed guarantee over a uniform signal, then apply
Theorem~\ref{s:main}.
\end{proof}

The dimension threshold is uniform in the accuracy; no relation between
accuracy and memory is assumed. For example, when
\(\epsilon=d^{-r}\le1/10\), the conclusion is
\(T\ge2^{-16}r d\log_2 d\).

\subsection*{Context and ancestry}


Memory restrictions affect statistical risk already in sparse noisy
regression, as shown by Steinhardt and Duchi
\cite[Section~1.1, equation~(2)]{SteinhardtDuchi2015}.
Steinhardt, Valiant, and Wager related bounded-memory inference to
communication and statistical-query models and posed a quadratic-memory
versus exponential-sample question for parity learning
\cite[Section~1.1, Conjecture~1]{SVW2016}. Raz established the
qualitative separation through finite-width branching programs
\cite[Theorem~1]{Raz2016}. These works make explicit the distinction
between information in a current observation and information that a
finite state can carry forward.

The closest continuous-regression comparison among these sources is
Sharan, Sidford, and Valiant \cite{SSV2019}. Their Theorem~1 uses
independent isotropic Gaussian covariates and independent uniform
additive noise of half-width \(2^{-d/5}\). At Euclidean accuracy
\(\epsilon=d^{-r}\), for \(r\le O(d/\log d)\) in the range specified
there, at most \(d^2/4\) bits and success probability at least \(2/3\)
require \(\Omega(d\log r)\) samples
\cite[Theorem~1]{SSV2019}. Their proof experiment draws the target
uniformly on the sphere \cite[Section~2]{SSV2019}. The exact labels in
Definition~\ref{s:model} are more informative than noisy labels, so
that noisy lower bound alone does not give Theorem~\ref{s:main}.
The condition-number conjecture in their Section~1.1 concerns a
different, ill-conditioned Gaussian experiment at constant accuracy.

Dagan, Kur, and Shamir give related space lower bounds for distinct
linear-prediction tasks \cite{DKS2019}. Their published Theorem~3 asks
for an approximate null vector of independent Gaussian rows. Their
Theorem~10 concerns empirical residual error on a finite equation
system formed from jointly conditioned rows and one fixed equation.
Those targets and observation laws differ from the fixed-signal
precision problem considered here.

Our block analysis uses high moments of several directions drawn in
one random subspace. The use of independent copies and successive
orthogonalization has a close predecessor in the high-moment
calculation of Sharan, Sidford, and Valiant
\cite[Section~7, Lemma~12]{SSV2019}. Their calculation concerns a noisy
interval projection under a global \(L^2\) density hypothesis. The
local argument below instead conditions an exact Gaussian block and
controls its residual subsphere. Classical beta laws for random
projections are discussed by Frankl and Maehara
\cite{FranklMaehara1990}; all normalized densities, conditional laws,
and inequalities used here are proved locally.

\subsection*{Exact observations and residual subspheres}

For a Euclidean linear space \(H\), write \(S(H)\) for its unit sphere
and \(\sigma_H\) for uniform probability on that sphere. We test the
success probability of a learner on affine subspheres
\begin{equation}\label{s:affine-sphere}
 z+rS(H),\qquad z\perp H,\quad r>0,\quad \norm z^2+r^2=1.
\end{equation}
The dimension of this test means the linear dimension \(D=\dim H\).
Fix the shared seed and the learner's rules. At a block boundary, let
\(g_v(s)\) be the success probability of the
remaining computation started in state \(v\), using fresh samples.
This function is defined even when \(v\) is unreachable at target
\(s\). The proof bounds its average on every test
\eqref{s:affine-sphere} of sufficiently large dimension.

To evaluate one such average, draw the target uniformly on the affine
test, independently of fresh block rows. A block of \(1\le b<D\)
Gaussian rows identifies the target's projection
\(p\) onto their row span in \(H\). Conditional on the full block
data, the target is uniform on
\[
 z+p+\rho S(K),\qquad \rho^2=r^2-\norm p^2,
\]
where \(K\) is the perpendicular complement of that row span in
\(H\). When the experiment is instead conditioned on the observed
projection \(p\), the space \(K\) is uniform among the
\((D-b)\)-dimensional subspaces of \(E=H\cap p^\perp\).
Thus the block leaves a random residual subsphere, even though the
learner may process its exact real data without computational limits.
Section~\ref{s:residual-spheres} proves these conditional laws.

The state chosen after the block may depend on its entire matrix and
all its labels. We account for this choice before averaging away the
matrix data: the resulting success probability is bounded by the
maximum over the possible continuation states. For fixed \(p,\rho\),
write
\[
 F_v(K)=\int g_v(z+p+\rho u)\,d\sigma_K(u).
\]
If the block can choose among at most \(W\) states, then for every
positive integer \(h\),
\[
 \E_K\max_v F_v(K)
 \le\left(\sum_v\E_K F_v(K)^h\right)^{1/h}.
\]
A bound on the \(h\)-th moment therefore charges state selection by
\(W^{1/h}\). The moment is an average over \(h\) directions drawn
independently in the same random space \(K\); those directions are
correlated after \(K\) is averaged out.

\subsection*{Proof of the precision bound}

The central task is to control this correlation without spending a
new test dimension for every direction. Conditional on preceding
directions with span \(V\), the next direction has more mass near
\(V\) than a uniform direction in \(E\). The comparison in
Section~\ref{s:beta-mixture} bounds its law by an exponential factor
times two simpler probability laws. The first is uniform on
\(S(E)\). The second chooses one direction in \(V\) and then draws a
uniform direction on a latitude perpendicular to it. Both laws can
be bounded by the same suffix estimate on affine subspheres: the
first uses dimension \(D-1\), and the second dimension \(D-2\).
Every factor in the moment uses these same tests. Thus a block costs
only two test dimensions, although it contains a number of samples
proportional to \(d\).

The residual radius and the comparison-latitude radii contribute
inverse-radius moments. Their squared radius fractions have beta
laws, and finite beta products bound these moments explicitly.
Together with a moment order proportional to \(d\), this gives a
uniform upper bound on success that can be propagated backward over
the blocks. The terminal bound is a spherical cap estimate.

There is still a dimension constraint: this induction is useful only
while the test spheres remain large. Section~\ref{s:stopping} supplies
that constraint from the learner model. There are at most
\((T+1)2^M\) terminal index-state pairs. Counting their cap
probabilities shows that a successful putative run with
\(T=O(d\log(1/\epsilon))\) must have
\(\log(1/\epsilon)=o(d)\) when \(M=o(d^2)\). Such a run contains
only \(o(d)\) blocks, so the induction has sufficient dimensions.
A full-data residual estimate handles bounded precision.

The direct proof of Theorem~\ref{s:main} therefore proceeds through
the continuation and residual laws in Section~\ref{s:residual-spheres},
the cap and common-subspace direction facts in
Section~\ref{s:subspace-moments}, the stopping reductions in
Section~\ref{s:stopping}, and the beta-mixture comparison in
Section~\ref{s:beta-mixture}.

\subsection*{Further analytic formulations}

The same geometric setup also gives bounds that apply to arbitrary
bounded continuation functions, whether or not they arise from a
learner. Section~\ref{s:weighted-proof} develops a second way to
control the correlated directions. A cap estimate bounds mass in a
tube around the preceding span, and summing shells controls the
singular conditional density. This yields
Theorem~\ref{s:weighted-theorem}: a block selecting among at most
\(W\) continuations increases their radius-weighted subsphere means
by at most \(W^{1/h}C^d\), with \(h=\lfloor d/16\rfloor\),
while lowering the test dimension by two.
Its terminal bound gives the same precision rate at success
probability \(1/2\), with an unspecified absolute coefficient.

Section~\ref{s:affine-invariant} gives an explicit version of the
suffix estimate simultaneously on all admissible affine subspheres.
An even radius exponent lets its inverse-radius moment be evaluated
by a finite beta product, retaining explicit constants throughout.
Section~\ref{s:fixed-length} gives a different extension: its
induction compares actual learners in different ambient dimensions
and permits arbitrary Euclidean outputs at every positive error
tolerance. The key step simulates a latitude target by a smaller
Gaussian learner without increasing its state set. These are
additional formulations of the method; the numerical coefficient in
Theorem~\ref{s:main} comes from the beta-mixture proof.
